\documentclass[11pt]{article}

\usepackage[margin=1in]{geometry}
\usepackage{amsmath,amssymb,amsthm}
\usepackage{booktabs}
\usepackage{mathtools}
\usepackage{enumitem}
\usepackage[dvipsnames]{xcolor}
\usepackage{microtype}
\usepackage[hidelinks]{hyperref}

\newtheorem{theorem}{Theorem}
\newtheorem{proposition}[theorem]{Proposition}
\newtheorem{lemma}[theorem]{Lemma}
\newtheorem{corollary}[theorem]{Corollary}
\theoremstyle{definition}
\newtheorem{definition}[theorem]{Definition}
\theoremstyle{remark}
\newtheorem{remark}[theorem]{Remark}

\newcommand{\F}{\mathrm{F}}
\newcommand{\GHZ}{\mathrm{GHZ}_3}
\newcommand{\LC}{\mathrm{LC}}
\newcommand{\rw}{\mathrm{rw}}
\newcommand{\Prod}{\mathsf{Prod}}
\newcommand{\N}{N}
\newcommand{\code}[1]{\texttt{#1}}

\title{\bf Distance-hereditary graphs characterize\\ minimum-fusion synthesis of graph states\\[4pt]
\large A machine-discovered result with a formally verified core}

\author{Sean Howell\\[2pt]
\small \texttt{sean.g.howell@gmail.com}}

\date{July 2026}

\begin{document}
\maketitle

\begin{abstract}
In fusion-based quantum computation a graph state $|G\rangle$ on $\N$ qubits is assembled from
small entangled resource states by destructive two-qubit fusion measurements. We study the
minimum number $\F(G)$ of fusions required in the model whose resource is the three-qubit
Greenberger--Horne--Zeilinger state and whose fusion is the type-II Bell measurement (the
$\GHZ$ model). Counting qubits and connected components gives the universal bound
$\F(G)\ge \N-3$ for connected targets. We determine exactly which graphs attain it:
$\F(G)=\N-3$ if and only if $G$ is \emph{distance-hereditary}---equivalently, of rank-width at
most one; equivalently, $C_5$-vertex-minor-free.

The result was produced by an autonomous formal-methods harness rather than by hand. A language
model (Fable~5), permitted to emit only machine-checkable artifacts and never given execution
access, classified all $585$ local-complementation orbits of connected graphs on at most nine
vertices by whether they attain the bound, every certificate accepted by two independent
stabilizer engines; proposed the distance-hereditary characterization from the $44$ labeled
orbits on at most seven vertices, after which all $541$ held-out orbits on eight and nine
vertices proved consistent; and formalized the core of the result in Lean~4 against an
adversarially hardened trust gate. Formalizing the forward direction
(distance-hereditary $\Rightarrow \F=\N-3$) required extending the production model with a second
fusion primitive, itself verified against both engines; the reverse direction is a Lean theorem in
a single-blob schedule model that over-approximates intra-blob fusions, where a formalized counting
argument shows that no schedule at the floor can use one, and the only step left to the physical
minimum is the reduction of an arbitrary schedule to a single growing blob by commutation of
fusions on disjoint qubit pairs. Under the harness's discipline no claim is recorded above the
status of a heuristic without a machine artifact, and the mathematics, the conjecture, and the
Lean proofs were authored by the model.
\end{abstract}

\section{Introduction}

Fusion-based quantum computation (FBQC) builds fault-tolerant computation from a supply of small,
identical entangled \emph{resource states} that are stitched together by \emph{fusions}:
destructive joint measurements of two qubits~\cite{fbqc}. Because the resource states and the
fusions are the physical cost drivers of a photonic architecture, a basic question for any target
entangled state is how few fusions suffice to prepare it. For graph states---the stabilizer states
in one-to-one correspondence with simple graphs~\cite{graphstates}---this is a combinatorial
optimization over the graph.

We work in a specific instance, which we call the $\GHZ$ model: the resource is the three-qubit
GHZ state $|\GHZ\rangle=\tfrac{1}{\sqrt2}(|000\rangle+|111\rangle)$, whose graph is the
three-vertex path, and each fusion is a type-II (destructive) Bell measurement. Since a stabilizer
state is defined only up to single-qubit Clifford operations, the natural cost is a function of
the local-Clifford orbit. Writing $\N=|V(G)|$, we set
\[
  \F(G)\;=\;\text{the least number of fusions producing a state locally Clifford equivalent to }|G\rangle,
\]
starting from a fresh supply of $|\GHZ\rangle$ states, with single-qubit Clifford operations free
throughout.

A short accounting of qubits and connected components gives a universal lower bound
$\F(G)\ge \N-3$ (Section~\ref{sec:model}). The graphs meeting it form the extremal class of the
problem, and our main result identifies that class exactly.

\begin{theorem}[Characterization]\label{thm:main}
For a connected graph $G$ on $\N\ge 3$ vertices in the $\GHZ$ model,
\[
  \F(G)=\N-3 \quad\Longleftrightarrow\quad G \text{ is distance-hereditary.}
\]
\end{theorem}

Distance-hereditary graphs~\cite{bandeltmulder,howorka} are exactly the graphs of rank-width at
most one~\cite{oum-rw}, the graphs with no vertex-minor isomorphic to the five-cycle $C_5$, and
the graphs obtainable from a single vertex by repeated addition of pendant vertices and twins.
Rank-width is invariant under local complementation~\cite{oum-rw,bouchet}, and $\F$ is invariant
under local Clifford operations; the characterization thus matches the fusion cost to a
local-Clifford-invariant measure of the entanglement structure of $|G\rangle$.

\paragraph{Contributions.} The mathematical content of the paper is Theorem~\ref{thm:main} and its
constituent parts:
\begin{enumerate}[nosep,leftmargin=1.4em]
\item a certified classification of all $585$ connected local-complementation orbits with
      $\N\le 9$ by whether $\F=\N-3$, together with exact values strictly above the bound for
      thirteen orbits; every certificate carries a witness accepted by two independent stabilizer
      engines (Section~\ref{sec:tablebase});
\item Lean~4 proofs of the universal lower bound and of $\F=\N-3$ for paths, stars, all trees,
      complete graphs $K_\N$, and complete bipartite graphs $K_{m,n}$
      (Section~\ref{sec:families});
\item a Lean~4 proof of the forward direction of Theorem~\ref{thm:main}
      (distance-hereditary $\Rightarrow \F=\N-3$), obtained by extending the production model with
      a second, engine-verified fusion primitive (Section~\ref{sec:forward}); and
\item a Lean~4 equivalence identifying the graphs producible at the universal floor with the
      distance-hereditary graphs, in a single-blob schedule model whose \code{intra} constructor
      over-approximates intra-blob fusions (\code{floor\_schedule\_iff\_dh}); the lift to the
      physical minimum now rests only on the single-blob normal form, a reordering justified by
      the commutation of fusions on disjoint qubit pairs (Section~\ref{sec:reverse}).
\end{enumerate}

The paper also documents the method. Every item above was discovered or written by an autonomous
harness driving a single language model under an explicit epistemic discipline: no statement is
recorded above the status \textsc{heuristic} without an accompanying machine artifact, the model
that does the mathematics never executes anything itself, and the formalization is checked by a
trust gate that was hardened against adversarial submissions. We describe the harness in
Section~\ref{sec:methods} because the provenance bears on how the results should be read.

\paragraph{What is proved, and in what sense.} Because these distinctions matter for a
machine-aided result, we state them plainly and keep to them throughout.
\begin{itemize}[nosep,leftmargin=1.4em]
\item \emph{Lean theorem}: a statement checked by the Lean~4 kernel and a hardened gate, using only
      the axioms \code{propext}, \code{Classical.choice}, \code{Quot.sound}, against
      \textsc{mathlib} pinned at \code{v4.31.0}. The lower bound, the five families, the forward
      direction, the model-level equivalence, and the counting bridge in the single-blob schedule
      model of Section~\ref{sec:reverse} are of this kind.
\item \emph{Engine-verified}: a fact about the physical fusion operation confirmed by two
      independently implemented stabilizer simulators that share no code. The two fusion primitives
      used by the production model are of this kind; inside Lean they are graph rewrites, and their
      right to be called ``one fusion'' rests on the engine agreement rather than on a derivation
      of qubit semantics in Lean.
\item \emph{Prose argument}: the single-blob normal form of Section~\ref{sec:reverse}, which shows
      that an arbitrary multi-component schedule can be reordered, by commutation of fusions on
      disjoint qubit pairs, into one that grows a single blob. It is proved here at the ordinary
      standard of rigor of a mathematics paper and rests on a domain-level, engine-checked
      reduction; it is not formalized. The counting bridge itself, previously in this category, is
      now a Lean theorem in the single-blob model (Section~\ref{sec:reverse}).
\end{itemize}
Theorem~\ref{thm:main} combines all three. Independently of its proof, the biconditional was
checked against the certified tablebase on all $585$ orbits and holds with no exception.

\section{The \texorpdfstring{$\GHZ$}{GHZ3} fusion model and the lower bound}\label{sec:model}

\paragraph{Graph states and local complementation.} To a simple graph $G=(V,E)$ one associates the
stabilizer state $|G\rangle=\prod_{\{u,v\}\in E}\mathrm{CZ}_{uv}\,|{+}\rangle^{\otimes V}$. Local
Clifford equivalence of graph states is generated by \emph{local complementation}: complementing at
a vertex $v$ toggles the adjacency of every pair of neighbors of $v$~\cite{vandennest,bouchet}. We
write $G\sim_{\LC}H$ for the resulting equivalence and note that both $\F$ and the rank-width of $G$
are constant on $\sim_{\LC}$ orbits.

\paragraph{Resource and fusion.} The resource $|\GHZ\rangle$ has graph the three-vertex path $P_3$.
A type-II fusion destructively measures the two-qubit Paulis $X\otimes Z$ and $Z\otimes X$ on a
pair of qubits; on a graph state it removes the two measured vertices and rewires their
neighborhoods, up to Pauli byproducts that are absorbed into the local-Clifford equivalence.
Because intermediate and final states are tracked only up to local Cliffords, fusion conventions
that differ by single-qubit Cliffords on the measured pair yield the same cost $\F$.

\paragraph{Model dependence.} The value $\F$ is specific to this resource and this fusion. Other
natural models---for instance models with emitter-generated caterpillar resources---assign
different values, and the extremal class can differ between models. We fix the $\GHZ$ model
throughout and use externally tabulated values from other models only as soft cross-checks, never
as ground truth.

\paragraph{The universal lower bound.} Suppose a schedule consumes $g$ copies of $|\GHZ\rangle$ and
performs $f$ fusions to produce a connected $\N$-vertex graph state. Each resource contributes
three qubits and each fusion consumes two, so
\begin{equation}\label{eq:qubits}
  3g-2f=\N.
\end{equation}
The $g$ resources begin as $g$ connected components and the target is a single component.
Components merge only when a fusion consumes one qubit from each of two distinct components, and a
resource that never merges into the final component must instead have all three of its qubits
consumed by fusions, which costs more than merging it; an induction over the schedule gives
$f\ge g-1$. Substituting $g\le f+1$ into~\eqref{eq:qubits} yields
$\N=3g-2f\le 3(f+1)-2f=f+3$, hence

\begin{proposition}[Universal lower bound]\label{prop:lb}
$\F(G)\ge \N-3$ for every connected $G$ on $\N$ vertices.
\end{proposition}

This is the content of the Lean declaration \code{fusion\_cost\_lower\_bound}. The component
bookkeeping is not assumed but derived from a merge model of the schedule; in particular the Lean
proof covers the degenerate steps in which a fusion splits a component or consumes one outright,
which the one-line argument above glosses over.

\section{The Empiricist harness}\label{sec:methods}

The results were obtained by a harness (``Empiricist'') that drives a language model, Fable~5,
through research tasks while forbidding it any unverified step. Four of its design commitments bear
on how the results should be read.

\paragraph{An append-only epistemic ledger.} Every artifact---a computed value, a construction, a
conjecture, a Lean module---enters a ledger at a status in the ordered set
\[
  \textsc{refuted}<\textsc{heuristic}<\textsc{conjectured}<\textsc{verified}_{\N}<\textsc{certified}<\textsc{formalized},
\]
and rises only when accompanied by an evidence row from a verifier. \textsc{refuted} is terminal.
Nothing the model merely asserts is recorded above \textsc{heuristic}. Artifacts are addressed by
content hash and certificates embed exact version pins, so results are reproducible from the ledger
alone.

\paragraph{The model never executes.} Fable~5 emits structured JSON---a conjecture, a construction,
or a complete Lean module---and the harness runs every program in a sandbox. This keeps the
boundary between a model's suggestion and a checked fact sharp, and it is what lets us attribute
the mathematics to the model while still calling the results verified. The harness itself---ledger,
executor, verifiers, gate---is conventional software.

\paragraph{Two independent physical verifiers.} Fusion outcomes are checked by two stabilizer
engines that share no code: a tableau simulator built on \textsc{stim}~\cite{stim} and a GF(2)
bitmask implementation written independently from the physical specification. A fusion move is
accepted only when both engines agree on the resulting stabilizer state up to local Clifford
equivalence; disagreement is treated as a fault that halts the campaign. The closed-form graph
rewrite for each resource merge is trusted only because it matches both engines on a fuzz test
against random input states, cross-checked against \textsc{nauty}'s canonical
forms~\cite{nauty}.

\paragraph{A hardened Lean trust gate.} A submitted Lean~4~\cite{lean4} module is elaborated in a
sandbox with network and process creation denied and an ephemeral write jail, against a frozen
read-only snapshot of a pinned \textsc{mathlib}~\cite{mathlib}. The kernel result is re-checked by
an independent checker; the axiom set is required to lie within
$\{\code{propext},\code{Classical.choice},\code{Quot.sound}\}$; every import must resolve to a
member of an explicit trusted set; and the build directory is swept for residue after elaboration.
The gate was developed against successive rounds of adversarial red-teaming (poisoned olean files,
time-of-check/time-of-use swaps, axiom-hiding elaborators), with each exploit repaired before the
gate was used in production.

\paragraph{The formalization loop.} To formalize a statement the model emits a complete Lean module
as JSON; the gate either accepts it or returns structured diagnostics, including the exact proof
obligations remaining at unfilled holes. The model revises across rounds until the module passes.
The loop uses no shell and no interactive proof assistant on the model's side; the model develops a
proof by reading goal states the gate reports, not by running Lean.

\section{The certified tablebase}\label{sec:tablebase}

The harness organized the connected graphs with $3\le\N\le 9$ into local-complementation orbits
(relabeling and local complementation identified); there are $585$. For each orbit the ledger
records whether $\F=\N-3$, with a machine certificate either way: when the orbit is extremal, a
fusion schedule with $\N-3$ fusions whose outcome both engines accept, meeting the lower bound of
Proposition~\ref{prop:lb}; when it is not, a certified refutation of every schedule at the floor.
(By the counting argument of Section~\ref{sec:reverse}, a schedule with $\N-3$ fusions consumes
$\N-2$ resources and every one of its fusions joins two components, so the floor schedules form a
finite space; each fusion outcome in the exhaustion is computed by both engines.) The resulting
object---$585$ rows, each an orbit representative with its verdict and certificates---is a single
content-addressed artifact at status $\textsc{verified}_\N$.

\begin{table}[t]
\centering
\begin{tabular}{lrrrrrrrr}
\toprule
$\N$ & 3 & 4 & 5 & 6 & 7 & 8 & 9 & total\\
\midrule
connected orbits           & 1 & 2 & 4 & 11 & 26 & 101 & 440 & 585\\
extremal ($\F=\N-3$)       & 1 & 2 & 3 &  8 & 15 &  42 & 104 & 175\\
non-extremal ($\F>\N-3$)   & 0 & 0 & 1 &  3 & 11 &  59 & 336 & 410\\
\bottomrule
\end{tabular}
\caption{Connected local-complementation orbits by vertex count, split by whether the orbit attains
the universal minimum $\F=\N-3$. The extremal set has $175$ orbits; all are distance-hereditary.}
\label{tab:counts}
\end{table}

Table~\ref{tab:counts} records the split. Two further outcomes of the campaign are worth noting.
First, exact values strictly above the floor were closed for thirteen orbits. Ten sit one tier up,
at $\F=\N$ (one orbit at $\N=5$, two at $\N=6$, seven at $\N=7$), and were closed by deterministic
search. Three sit two tiers up, at $\F=\N+3$, and were closed by model-proposed schedules: a
six-vertex three-regular orbit at $\F=9$, realized with eight resources, nine fusions, and three
local-complementation steps, and two seven-vertex orbits at $\F=10$; each schedule was
independently re-verified by both engines. Second, the natural refinement
$\F=\N-3+3(\rw(G)-1)$, which would make the fusion cost a linear function of rank-width, is false:
a seven-vertex orbit of rank-width two carries a certified bound $\F\ge \N+3=10$, whereas the
formula predicts $\F=\N=7$. Rank-width captures the extremal boundary exactly but does not
determine the cost above it.

\section{The extremal class is distance-hereditary}\label{sec:characterization}

\paragraph{Discovery.} The harness presented the model with the complete labeled catalogue of the
$44$ orbits with $\N\le 7$, each as an edge list tagged extremal or not, with no invariant, class
name, or hint in the prompt beyond the true observation that the target property is a
local-complementation invariant. The model proposed that the extremal orbits are exactly the
distance-hereditary graphs, supplying the standard equivalent characterizations (metric, rank-width
at most one, forbidden induced house/gem/domino/long-cycle, and pendant/twin reducibility), the
relevant invariance results~\cite{bandeltmulder,oum-rw,bouchet}, and the identification of the
smallest obstruction: the unique non-extremal five-vertex orbit is the local-complementation orbit
of $C_5$.

\paragraph{Grounding.} A deterministic checker---Bandelt--Mulder pendant/twin elimination, which
reduces a graph to a single vertex if and only if it is distance-hereditary---was run on all $585$
orbits and compared against the extremal label. The classification $\F=\N-3\Leftrightarrow$
distance-hereditary holds with zero exceptions, including on the $541$ held-out orbits with
$\N\in\{8,9\}$ that the model never saw. The proposal was recorded as a \textsc{conjectured}
artifact, promoted from \textsc{heuristic} by this grounding.

Every family for which we prove $\F=\N-3$ in Section~\ref{sec:families} is distance-hereditary,
which is why each attains the bound: trees, complete graphs, and complete bipartite graphs all
have rank-width at most one.

\section{Formalized families and the lower bound}\label{sec:families}

The Lean development begins with Proposition~\ref{prop:lb}, formalized as a self-contained module,
and with a production predicate. $\Prod_k(G)$ holds when the concrete graph $G$ is obtained from
$P_3$ by $k$ successive \emph{leaf-merges}, up to isomorphism; a leaf-merge attaches one pendant
vertex, and its graph rewrite is proved in Lean to be isomorphic to that of the physical resource
merge at a leaf qubit, whose physical meaning is engine-verified (Section~\ref{sec:methods}). The
predicate deliberately has no local-complementation constructor, so it generates exactly the
trees; local equivalence enters only when a statement about $\F$ is made.

\begin{theorem}[Families]\label{thm:families}
In the $\GHZ$ model, $\F(G)=\N-3$ for $G$ a path, a star, any tree, a complete graph $K_\N$, or a
complete bipartite graph $K_{m,n}$. Each is a Lean theorem: the lower bound is
Proposition~\ref{prop:lb}, and the upper bound exhibits a graph $H\sim_{\LC}G$ with
$\Prod_{\N-3}(H)$.
\end{theorem}

The proofs reduce each family to a producible tree up to local complementation: $K_\N$ is one local
complementation from a star, $K_{m,n}$ is three from a double star, and trees are produced directly
by leaf-removal induction, of which paths and stars are special cases. The statements are faithful
in a strict sense: each carries the general lower bound in its own statement---quantifying over
arbitrary fusion schedules, not only the exhibited construction---so that a theorem cannot be true
for a vacuous or circular reason.

\begin{remark}[Induction must track concrete graphs]\label{rem:naive}
Although every statement about $\F$ is invariant under local complementation, the inductions cannot
be conducted on $\sim_{\LC}$ orbits, because pendant addition does not descend to orbits:
$P_3\sim_{\LC}K_3$, yet attaching a pendant to the center of $P_3$ yields the four-vertex star,
while attaching a pendant to a vertex of $K_3$ yields the paw, which is LC-equivalent to $P_4$ and
not to the star. The Lean development therefore carries exact graphs through every induction and
applies local complementation only explicitly.
\end{remark}

\section{The forward direction of the characterization}\label{sec:forward}

Theorem~\ref{thm:families} covers exactly the graphs locally equivalent to a tree: since $\Prod$
generates the trees and nothing else, it can witness the floor only on the local-complementation
orbit of a tree. This is a strict subclass of the distance-hereditary graphs. A direct check on
the tablebase shows that of the $175$ extremal orbits, $93$ contain a tree in their
local-complementation orbit and $82$ do not---two at $\N=6$, four at $\N=7$, nineteen at $\N=8$,
and fifty-seven at $\N=9$. (The orbit enumerations behind these counts ran to completion; a safety
cap of $200000$ graphs per orbit was never reached.) The forward direction of
Theorem~\ref{thm:main} is therefore not provable in that model.

\paragraph{The missing primitive.} The obstruction is a missing operation, not a difficult proof.
The physical resource merge fuses one qubit $a$ of the current state with one qubit of a fresh
$|\GHZ\rangle$, and up to the symmetry of $P_3$ there are two choices: the fused resource qubit is
a leaf or the center. The leaf choice attaches a pendant. The center choice removes the fused
state vertex $a$ and introduces the two resource leaves, each adjacent to exactly the former
neighborhood of $a$ and not to each other; up to relabeling this is $\mathrm{addFalseTwin}(G,a)$.
In Lean the center-merge rewrite is proved isomorphic to false-twin addition
(\code{ghz3CenterMerge\_iso\_addFalseTwin}); physically, the closed form matches the engine
outcome up to local Clifford equivalence on $400$ of $400$ random input states, with the two
engines agreeing with each other on all $400$---the same discipline already applied to the
leaf-merge. The center merge is a genuine one-fusion primitive; the production model had simply
never included it.

\paragraph{The extended model.} We formalize the production predicate $\Prod^{+}_k$ generated by
five constructors: the base $P_3$; the leaf-merge (pendant) and center-merge (false twin), each
adding one vertex and one fusion; local complementation, which is free; and graph isomorphism.
Because distance-hereditary graphs are closed under pendant addition, twin addition, local
complementation, and isomorphism, every $\Prod^{+}$-producible graph is distance-hereditary. The
converse is the content of the forward direction.

\begin{theorem}[Forward direction, formalized]\label{thm:forward}
Let $\mathsf{DH}$ denote the class generated from the connected three-vertex graphs ($P_3$ and
$K_3$) by pendant, false-twin, and true-twin additions, closed under isomorphism; these are
exactly the connected distance-hereditary graphs on at least three
vertices~\cite{bandeltmulder}. Then, as a Lean theorem,
\[
  G\in\mathsf{DH}\ \Longrightarrow\ \Prod^{+}_{\N-3}(G),
\]
and hence, with Proposition~\ref{prop:lb}, $\F(G)=\N-3$.
\end{theorem}

The proof is a structural induction on the Bandelt--Mulder build, carrying the exact graph as
Remark~\ref{rem:naive} requires. A pendant addition maps to a leaf-merge; a false-twin addition
maps to a center-merge; a true-twin addition is realized as local complementation, then a
leaf-merge, then local complementation, using the graph identity
\[
  \mathrm{lc}\big(\mathrm{addPendant}(G,u),\,u\big)\;=\;\mathrm{addTrueTwin}\big(\mathrm{lc}(G,u),\,u\big)
\]
together with the involutivity of local complementation. No search for a witness is needed:
tracking the exact graph through the extended model turns each build step into one constructor.
This is what makes the theorem short once the extended model is in place, and it is the step at
which the ``locally equivalent to a tree'' obstruction dissolves.

\section{The reverse direction}\label{sec:reverse}

The reverse direction is proved at the level of the extended model and then lifted to the physical
minimum by counting.

\begin{theorem}[Model characterization, formalized]\label{thm:reverse}
As a Lean theorem,
\[
  \Prod^{+}_{\N-3}(G)\ \Longleftrightarrow\ G\in\mathsf{DH}.
\]
\end{theorem}

The right-to-left implication is Theorem~\ref{thm:forward}. The left-to-right implication is a
structural induction on $\Prod^{+}$: the base is a three-vertex core, the leaf-merge and
center-merge map back to a pendant and a false twin through the same faithfulness isomorphisms
used above, and local complementation and isomorphism map to themselves. Making the last two steps
immediate requires recording that $\mathsf{DH}$ is closed under local complementation; since
rank-width at most one is a local-complementation invariant this does not enlarge the class, and
it replaces what would otherwise be a family of transport identities by a single constructor.

\paragraph{From the model to the physical minimum.} Theorem~\ref{thm:reverse} concerns the extended
model, in which every fusion attaches a fresh resource. The gap between that model and the physical
minimum has two parts: a schedule may perform \emph{intra-blob} fusions, consuming two qubits of
the same component, and it may grow several components in parallel before merging them. The first
part is now internalized in Lean. We formalize a schedule model $\mathsf{BlobSchedule}_f(G)$ (Lean
inductive \code{BlobSchedule}, module \code{Bridge}) with the five constructors of $\Prod^{+}$
together with a sixth, \code{intra}, which over-approximates an intra-blob fusion: from a graph on
a carrier $V$ it may produce \emph{any} graph on a carrier whose cardinality is smaller by two, at
the cost of one fusion, with no rewrite assumed. Whatever an intra-blob fusion does physically---it
removes the two measured qubits and rewires the rest in a way the harness computes only through the
engines---it is an instance of this constructor. Two counting facts then close the case in Lean.
The lemma \code{BlobSchedule.natCard\_le} shows that any $\mathsf{BlobSchedule}_f(G)$ has at most
$f+3$ vertices: each merge gains one vertex per fusion, and an intra-fusion loses two vertices
while still spending a fusion. The theorem \code{blobSchedule\_floor} shows that at the floor
$|V|=f+3$ the \code{intra} case is impossible---its sub-derivation would violate the lemma---so a
floor schedule collapses into $\Prod^{+}_f$. Combined with Theorem~\ref{thm:reverse} this yields:

\begin{theorem}[Counting bridge, formalized]\label{thm:bridge}
As a Lean theorem (\code{floor\_schedule\_iff\_dh}), for $\N\ge 3$,
\[
  \mathsf{BlobSchedule}_{\N-3}(G)\ \Longleftrightarrow\ G\in\mathsf{DH}.
\]
\end{theorem}

Even a model generous enough to let an intra-fusion produce an arbitrary graph on two fewer
vertices cannot use one at the floor, and the floor is attained exactly by the distance-hereditary
graphs.

The second part---the \emph{single-blob normal form}---is the one step that remains a prose
argument. Suppose a physical schedule attains $f=\F(G)=\N-3$ using $g$ resources. Since a fusion
destroys both of its qubits, the fusions of a schedule act on pairwise disjoint qubit pairs and
commute. Reordering the inter-component fusions along a root-first traversal of the merge structure
on the resources, and placing each intra-component fusion after the resources carrying its qubits
have been attached, gives an equivalent schedule in which every fusion either merges one fresh
resource into a single growing blob or acts inside that blob. Up to the free local Cliffords a
fresh resource carries the graph $P_3$ or its local complement $K_3$, so each merge fuses the
growing state with a leaf of $P_3$, the center of $P_3$, or a vertex of $K_3$; on graphs these
realize a pendant, a false twin, or a true twin respectively, all available in the model (the true
twin via the identity of Section~\ref{sec:forward}), local Cliffords on the growing state are free
local complementations, and each intra-blob fusion is an instance of \code{intra}. The reordered
schedule is therefore a $\mathsf{BlobSchedule}_{\N-3}$ derivation, and Theorem~\ref{thm:bridge}
gives that $G$ is distance-hereditary. (At the floor, \eqref{eq:qubits} gives $g=\N-2$ and the
component bound holds with equality, so in fact every fusion joins two components; but the argument
no longer depends on that exclusion, which now lives inside Lean as \code{blobSchedule\_floor}.)
Combined with Theorem~\ref{thm:forward} this proves Theorem~\ref{thm:main}.

\begin{remark}
The previous version of this paper reported the counting bridge as a prose argument, not
internalized as a single Lean theorem, on the ground that an intra-component fusion has no
closed-form graph rewrite---the harness computes it only through the stabilizer engines---so no
Lean model of schedules with intra-fusions seemed available. The over-approximation in
\code{BlobSchedule} sidesteps that obstacle: excluding an intra-fusion at the floor requires no
rewrite at all, only the vertex count, which is exactly what \code{BlobSchedule.natCard\_le} and
\code{blobSchedule\_floor} record. What remains outside Lean is only the single-blob normal
form---the reduction of an arbitrary multi-component schedule to a single growing blob, justified
by the commutation of fusions on disjoint qubit pairs; it is a domain-level reduction, checked by
the engines, and not formalized. The alternative classical route to the reverse direction (every
graph that is not distance-hereditary has $\F>\N-3$) would still run through rank-width lower
bounds~\cite{oum-rw,oum-seymour}, and rank-width is not present in \textsc{mathlib}; the
over-approximation makes that route unnecessary for this step.
\end{remark}

\section{Discussion}

\paragraph{The modeled boundary.} The two fusion primitives are the one place where the formal
development rests on something other than the Lean kernel. Inside Lean they are graph rewrites;
their right to be called ``one fusion'' comes from agreement between two independent stabilizer
engines and a canonical-form cross-check, not from a derivation of qubit semantics in Lean. This is
a deliberate boundary, stated rather than hidden: the combinatorial optimization is proved, and the
physical primitive it optimizes over is pinned down by redundant simulation. A fully
first-principles account would formalize the stabilizer formalism and the Bell measurement in Lean
and re-derive the two rewrites; that is a separate and larger undertaking.

\paragraph{Discovery under discipline.} The characterization was proposed by the model from labeled
data and then survived an exhaustive deterministic check; the forward direction was formalized by
the model against a gate it could not talk its way past. The discipline has a measurable payoff
beyond trust: the deterministic recheck of the characterization corrected an arithmetic error in an
earlier hand-written note (the extremal/non-extremal split is $175/410$, not $185/400$), and the
faithfulness requirement on Theorem~\ref{thm:families} caught a near-circular statement in which
the construction class had been silently baked into the claim. A workflow that permits unchecked
assertion would have surfaced neither.

\paragraph{Open questions.} Three stand out. First, the remaining step of the physical reverse
direction: the single-blob normal form. The obstacle identified in the previous version of this
paper---that an intra-component fusion has no closed-form rewrite---is gone, sidestepped by the
over-approximation of Section~\ref{sec:reverse}; what remains open is to formalize the reduction of
an arbitrary multi-component schedule to a single growing blob, that is, a Lean model of general
schedules together with the reordering argument by commutation of fusions on disjoint qubit pairs.
Second, a finer invariant governing $\F$ above the extremal boundary, where rank-width is
insufficient (Section~\ref{sec:tablebase}). Third, the complexity of computing $\F$: the
characterization makes the decision problem $\F(G)\stackrel{?}{=}\N-3$ solvable in polynomial
time, since distance-hereditary graphs are recognizable in polynomial time via pendant/twin
elimination~\cite{bandeltmulder}, but the complexity of computing $\F$ in general is open.

\section{Conclusion}

In the $\GHZ$ fusion model the graphs that attain the universal minimum fusion count are precisely
the distance-hereditary graphs. The characterization was proposed by a language model from
certified data, holds with no exception on all $585$ connected local-complementation orbits with at
most nine vertices, and is formalized in Lean~4 in the forward direction and, in the reverse
direction, as a floor theorem for a single-blob schedule model that over-approximates intra-blob
fusions; the one step connecting that theorem to the physical minimum---the reduction of an
arbitrary schedule to a single growing blob---remains a prose argument backed by the engines. The
development was carried out end to end by an autonomous harness under a discipline that records
nothing above a heuristic without a machine artifact; that discipline is both how the result was
obtained and the sense in which it should be trusted.

\appendix
\section{Principal Lean declarations}

Table~\ref{tab:lean} lists the main formalized statements. All are checked by the hardened gate and
use only the axioms $\{\code{propext},\code{Classical.choice},\code{Quot.sound}\}$. The foundation
modules are committed and reusable; the theorem modules are recorded as gate-certified ledger
artifacts.

\begin{table}[h]
\centering
\small
\begin{tabular}{lll}
\toprule
declaration & statement & role\\
\midrule
\code{fusion\_cost\_lower\_bound} & $\F(G)\ge \N-3$ (schedule form) & Prop.~\ref{prop:lb}\\
\code{pathGraph\_min\_fusions}    & $\F(\text{path}_\N)=\N-3$        & Thm.~\ref{thm:families}\\
\code{star\_min\_fusions}         & $\F(\text{star}_\N)=\N-3$        & Thm.~\ref{thm:families}\\
\code{tree\_min\_fusions}         & $\F(T)=\N-3$ for all trees       & Thm.~\ref{thm:families}\\
\code{complete\_min\_fusions}     & $\F(K_\N)=\N-3$                  & Thm.~\ref{thm:families}\\
\code{completeBipartite\_min\_fusions} & $\F(K_{m,n})=\N-3$          & Thm.~\ref{thm:families}\\
\code{ghz3CenterMerge\_iso\_addFalseTwin} & center-merge $=$ false twin & Sec.~\ref{sec:forward}\\
\code{dh\_forward} / \code{dh\_min\_fusions} & $\mathsf{DH}\Rightarrow \Prod^{+}_{\N-3}$ & Thm.~\ref{thm:forward}\\
\code{dh\_characterization}       & $\Prod^{+}_{\N-3}(G)\Leftrightarrow G\in\mathsf{DH}$ & Thm.~\ref{thm:reverse}\\
\code{floor\_schedule\_iff\_dh}   & $\mathsf{BlobSchedule}_{\N-3}(G)\Leftrightarrow G\in\mathsf{DH}$ & Thm.~\ref{thm:bridge}\\
\bottomrule
\end{tabular}
\caption{Principal formalized declarations, against \textsc{mathlib} pinned at \code{v4.31.0}.}
\label{tab:lean}
\end{table}

\begin{thebibliography}{10}

\bibitem{fbqc}
S.~Bartolucci, P.~Birchall, H.~Bomb\'in, et al.,
\newblock Fusion-based quantum computation,
\newblock \emph{Nature Communications} \textbf{14}, 912 (2023).

\bibitem{graphstates}
M.~Hein, J.~Eisert, and H.~J.~Briegel,
\newblock Multiparty entanglement in graph states,
\newblock \emph{Physical Review A} \textbf{69}, 062311 (2004).

\bibitem{vandennest}
M.~Van den Nest, J.~Dehaene, and B.~De Moor,
\newblock Graphical description of the action of local Clifford transformations on graph states,
\newblock \emph{Physical Review A} \textbf{69}, 022316 (2004).

\bibitem{bouchet}
A.~Bouchet,
\newblock Graphic presentations of isotropic systems,
\newblock \emph{Journal of Combinatorial Theory, Series B} \textbf{45}, 58--76 (1988).

\bibitem{bandeltmulder}
H.-J.~Bandelt and H.~M.~Mulder,
\newblock Distance-hereditary graphs,
\newblock \emph{Journal of Combinatorial Theory, Series B} \textbf{41}, 182--208 (1986).

\bibitem{howorka}
E.~Howorka,
\newblock A characterization of distance-hereditary graphs,
\newblock \emph{The Quarterly Journal of Mathematics} \textbf{28}, 417--420 (1977).

\bibitem{oum-rw}
S.~Oum,
\newblock Rank-width and vertex-minors,
\newblock \emph{Journal of Combinatorial Theory, Series B} \textbf{95}, 79--100 (2005).

\bibitem{oum-seymour}
S.~Oum and P.~Seymour,
\newblock Approximating clique-width and branch-width,
\newblock \emph{Journal of Combinatorial Theory, Series B} \textbf{96}, 514--528 (2006).

\bibitem{lean4}
L.~de Moura and S.~Ullrich,
\newblock The Lean~4 theorem prover and programming language,
\newblock in \emph{Automated Deduction -- CADE 28}, Springer, 625--635 (2021).

\bibitem{mathlib}
The mathlib Community,
\newblock The Lean mathematical library,
\newblock in \emph{Proc.\ 9th ACM SIGPLAN Int.\ Conf.\ on Certified Programs and Proofs (CPP)}, 367--381 (2020).

\bibitem{stim}
C.~Gidney,
\newblock Stim: a fast stabilizer circuit simulator,
\newblock \emph{Quantum} \textbf{5}, 497 (2021).

\bibitem{nauty}
B.~D.~McKay and A.~Piperno,
\newblock Practical graph isomorphism, II,
\newblock \emph{Journal of Symbolic Computation} \textbf{60}, 94--112 (2014).

\end{thebibliography}

\end{document}